% ======================================================================
% Section 6 — Discussion
% Compatible with \documentclass[pdflatex,sn-mathphys-num]{sn-jnl}
% Bind to Stage 47 Results labels; no new numbers beyond Results
% ======================================================================
\section{Discussion}\label{sec:discussion}

The results separate two design questions that are often conflated in
graph-based traffic forecasting: \emph{which graph is supplied to the
backbone}, and \emph{how that graph is used over the input window}. The
discussion below reads the empirical pattern under that separation. We keep
interpretation at the level of statistical dependency and architecture; we
do not claim causal traffic mechanisms.

\subsection{Why Dense Graphs Fail and What Learned Structure Adds}\label{sec:disc-dense}

The physical road-network graph is a harmful default for these backbones at
the operating points we study. On Los-loop the physical adjacency is dense
(order $10^{3}$ entries including self-loops for $207$ sensors), yet the
graph-free identity control reduces RMSE by roughly a quarter to a third
relative to that graph at every horizon and on both datasets
(Sects.~\ref{sec:res-physical} and \ref{tab:tgcn-main}). A simple reading is
\emph{oversmoothing} or, more generally, harmful aggregation: when almost
every sensor is treated as a neighbor of every other, the graph convolution
mixes signals that are not jointly informative for forecasting, and the
recurrent or non-recurrent backbone cannot recover node-specific dynamics
from the smeared representation. Identity adjacency freezes that mixing
without removing the temporal path, which is consistent with the large,
seed-stable gap between Physical and NoSpatial.

Learned single graphs (GSL, cGSL) occupy an intermediate position. They
shrink the physical graph to a much smaller support (e.g.\ $28$
contemporaneous edges on Los-loop) and usually improve on Physical, but they
still lose to NoSpatial in every dataset$\times$horizon cell
(Sect.~\ref{sec:res-single}). Sparsity alone is not sufficient: at a matched
$30$-edge budget, random and top-correlation graphs are worse than no graph,
while DAGMA-placed multi-lag edges with per-lag consumption are better
(Sect.~\ref{sec:res-controls}, Los-loop PH1 only). What the single learned
graphs appear to add, relative to Physical, is a less harmful neighborhood;
what they fail to add, relative to NoSpatial, is enough useful structure to
pay for any residual mis-specification. The multi-lag experiments ask
whether that residual can be recovered when dependency structure is allowed
to be lag-specific and is consumed in a lag-aligned way.

This reading is consistent with the proximity-versus-dependency motivation:
physical connectivity is an imperfect proxy for functional co-movement, and
a dense proxy can be worse than no proxy. It does not imply that
graph-based forecasting is generally harmful; it implies that the default
reference graph should be tested against a graph-free control before
learned structure is credited with gains.

\subsection{Structure Versus Consumption}\label{sec:disc-consumption}

The central architectural observation is the dissociation in
Sect.~\ref{sec:res-dissociation}. The three lag-specific DAGMA graphs are
fitted once and shared. Consumed as a static union by GCN-MultiGSL, they
yield $9.78\pm 0.14$ RMSE on Los-loop PH1; consumed per timestep by
T-GCN-MultiGSL, the same artifacts yield $4.84\pm 0.11$. The union retains
$28$ of the $30$ lagged edge slots, so the bulk of the edge content is
present in both conditions. The opposite outcomes therefore \emph{cannot be
explained by the edge set alone}. Because GCN and T-GCN also differ in
backbone (non-recurrent versus recurrent), this comparison is not a fully
controlled isolation of consumption; it provides evidence that graph
consumption is a major determinant of whether the learned multi-lag
structure is useful.

Within the T-GCN family, the consumption ladder is equally informative.
Fixed cyclic lag assignment already beats NoSpatial on Los-loop
($7.8\%$ relative RMSE reduction at PH1). Adding three global scalars
(Weighted) changes almost nothing (differences $\le 0.01$ RMSE), so a single
static mixture over lags is not what drives the gain. Replacing that global
mixture with a per-node, per-timestep gate (Mix) adds a further
$0.35$--$0.40$ RMSE on Los-loop in $5/5$ seeds at every horizon
(Sect.~\ref{sec:res-mix}), without modifying the learned edge sets. The
results thus indicate that the utility of the learned multi-lag structure
depends critically on how the graphs are consumed: lag-aligned use helps
where the graphs are informative, and finer selection helps further. The
gate varies \emph{use}, not \emph{structure}; the graphs remain static
fitted artifacts.

The cGSL--GSL contrast is a miniature of the same theme. Binary
symmetrization of the same contemporaneous artifact helps the non-recurrent
GCN substantially (cGSL $5.76$ vs GSL $7.83$ on Los-loop PH1) and is nearly
irrelevant for T-GCN ($5.82$ vs $5.86$). One consistent explanation is
\emph{aggregation compatibility}: a symmetric aggregator is better matched
to a symmetric operator, whereas the recurrent cell can absorb directed
asymmetry without the same penalty. We offer this as an observation of the
current study, not as a general theorem about DAG versus cyclic graphs, and
we do not revive a temporal reading of the contemporaneous DAG.

\subsection{What the Learned Graph Represents}\label{sec:disc-interpretation}

We describe the learned graphs as descriptive statistical dependency
structure under a linear DAGMA fit with an acyclicity regularizer. The
multi-lag construction stacks explicit lag blocks, and the consumption
ablations show that treating those blocks as lag-specific operators is more
effective than merging them into one static adjacency. That pattern is
\emph{consistent with} a lag-structured reading of the fitted matrices, but
we did not independently validate the lag interpretation against ground-truth
propagation, and acyclicity is an optimization constraint rather than causal
evidence.

The dataset boundary sharpens this reading. On SZ-Taxi the multi-lag fit
retains only two lagged edges, and every learned variant remains within seed
variability of the graph-free control (Sect.~\ref{sec:res-boundaries}). On
Los-loop, where the lag blocks are substantially denser
($12/3/15$ slots; union $28$), lag-aligned consumption yields the study's
positive result, and the $15$-minute resolution variant preserves a large
relative improvement over NoSpatial at $15$--$60$ minutes ahead
(Table~\ref{tab:los15}). The honest synthesis is conditional: learned
multi-lag graphs help forecasting when the fitted dependency structure is
informative enough to be worth consuming, and when the consumption mechanism
respects the lag organization of that structure. Where the fit is nearly
empty, the method collapses to the graph-free baseline.

None of this implies that the physical graph is useless in absolute terms,
or that learned graphs will help every network. It implies that the default
practice of treating a dense proximity graph as ground truth, and of
crediting any learned graph that improves on it, is incomplete without a
graph-free control and without an explicit account of how the graph is
consumed over time.
